%%% PART — version avec retour à la ligne du titre
% --- soul : réglages & robustesse dans les signets PDF
\usepackage{soul}
\setul{2pt}{0.25ex}
\setulcolor{background2}

% Neutraliser \ul dans les signets hyperref
\pdfstringdefDisableCommands{%
  \def\ul#1{#1}%
}

% Robuste pour titres/ToC

% Soulignement sous tout le bloc (multi-lignes ok)
\newcommand{\underlinetitlePart}[2][\linewidth]{%
  \tikz[baseline=(text.south)]{
    \node[
      inner sep=0pt, text=black,
      text width=#1, align=left
    ] (text) {#2};
  }%
}
\newlength{\PartAvail}

\newcommand{\taggedparttitle}[1]{%
  % H1 pour la structure (à adapter si tu veux part=H1, section=H2, etc.)
  \tagstructbegin{tag=H1}%
  \tagmcbegin{tag=Span}%
  #1%
  \tagmcend
  \tagstructend
}

\newcommand{\parttitlewithbox}[1]{%
  \begingroup
  \phantom{\taggedparttitle{#1}}%
  % -- dimensions
  \pgfmathsetlengthmacro{\RectW}{0.5em}  % largeur du pavé à droite
  \pgfmathsetlengthmacro{\PadR}{0.6em}   % espace entre texte et pavé

  % largeur dispo pour le texte = \linewidth - (pavé + marge)
  \setlength{\PartAvail}{\dimexpr\linewidth-\RectW-\PadR\relax}
  

  \begin{tikzpicture}[baseline=(t.base)]
    % --- texte : \ul gère le soulignement multi-lignes ---
    \node[anchor=west, inner sep=0pt, text width=\PartAvail, align=left] (t)
      {\color{black}\MakeUppercase{\underlinetitle{#1}}};
    % --- pavé à droite : même hauteur que le texte ---
\path[
  fill=primarycolor,
  rounded corners=0.25em % ← ajuste ici
]
  ($(t.north east)+(\PadR,0)$)
    rectangle
  ($(t.south east)+(\PadR+\RectW,0)$);
  \end{tikzpicture}%
  \endgroup
}

% Formatage de \part (saut de page ok pour \part et \part*)

% --- retirer le saut de page automatique de \part
\usepackage{titlesec}
\titleclass{\part}{straight} % <- au lieu de "page"

% flag "première part ?"
\newif\iffirstpart
\firstparttrue

% --- MODIFICATION : Gestion du saut de page AVANT la commande \part ---
% Nécessite le package etoolbox (déjà chargé dans votre préambule)
\preto{\part}{%
  \iffirstpart
    \global\firstpartfalse % pas de saut pour la toute première partie
  \else
    \clearpage             % saut de page forcé avant que l'ancre ne soit créée
  \fi
}

\titleformat{\part}[block]
  {\robotoCondensed\bfseries\fontsize{26}{30}\selectfont}
  {}{0pt}{%
    % Plus de \clearpage ici, c'est géré par le \preto ci-dessus
    \PlaceAnchor
    \parttitlewithbox
  }
\titlespacing*{\part}{1cm}{0cm}{0cm}